% Formulario TEEA 4MAT — Argomento 2: l'impedenza
% Per rigenerare il PDF (dalla cartella strumenti, due volte per il numero di pagine):
%   pdflatex -output-directory=../pdf -jobname=02-formulario-impedenza formulario.tex
% ---------------------------------------------------------------------------
%  Stile comune dei formulari TEEA (pdfLaTeX, solo pacchetti standard
%  presenti in TeX Live e MiKTeX). Testo in Computer Modern Sans.
% ---------------------------------------------------------------------------
\documentclass[11pt,a4paper]{article}
\usepackage[a4paper, top=14mm, bottom=16mm, left=14mm, right=14mm, headheight=0pt, footskip=8mm]{geometry}
\usepackage[utf8]{inputenc}
\usepackage{amsmath, amssymb}
\renewcommand{\familydefault}{\sfdefault}
\newcommand{\ohm}{\ensuremath{\Omega}}
\newcommand{\uA}{\ensuremath{\mu}}
\newcommand{\per}{\ensuremath{\cdot}}
\usepackage{xcolor, tikz, multicol, booktabs, array, fancyhdr, lastpage, tabularx}
\usepackage[most]{tcolorbox}
\usetikzlibrary{arrows.meta}

\definecolor{verde}{HTML}{159957}
\definecolor{blu}{HTML}{155799}
\definecolor{inchiostro}{HTML}{1B2936}
\definecolor{tenue}{HTML}{606C71}
\definecolor{fondo}{HTML}{EEF3F0}
\definecolor{tensione}{HTML}{2152C0}
\definecolor{corrente}{HTML}{B3205F}
\definecolor{resistenza}{HTML}{0B7A62}
\color{inchiostro}

\setlength{\parindent}{0pt}
\setlength{\columnsep}{7mm}
\renewcommand{\arraystretch}{1.25}

% unità di misura in tondo e virgola decimale
\newcommand{\U}[1]{\,\mathrm{#1}}
\newcommand{\virg}{{,}}

% riquadro di un argomento
\newtcolorbox{blocco}[1]{enhanced, breakable=false, colback=white, colframe=verde, boxrule=0.6pt,
  leftrule=3pt, arc=1.5pt, left=6pt, right=6pt, top=4pt, bottom=5pt,
  title={\bfseries #1}, coltitle=verde, colbacktitle=fondo, fonttitle=\normalsize,
  toptitle=1pt, bottomtitle=1pt, before skip=0pt, after skip=4.5mm}

% formula grande centrata
\newcommand{\F}[1]{\par\smallskip{\centering\Large$\displaystyle #1$\par}\smallskip}
% legenda dei simboli
\newcommand{\leg}[1]{{\small\color{tenue}#1}\par}

% intestazione del formulario
\newcommand{\intestazione}[3]{%
  {\color{verde}\bfseries\small #1}\par\vspace{1mm}
  {\fontsize{19}{23}\selectfont\bfseries #2}\par\vspace{1mm}
  {\small\color{tenue} #3}\par\vspace{2mm}
  {\color{verde}\rule{\linewidth}{1.6pt}}\par\vspace{4mm}}

\pagestyle{fancy}
\fancyhf{}
\renewcommand{\headrulewidth}{0pt}
\fancyfoot[L]{\scriptsize\color{tenue}Prof. Pasquale Naclerio \per{} TEEA 4MAT \per{} Formulario}
\fancyfoot[R]{\scriptsize\color{tenue}Pagina \thepage\ di \pageref{LastPage}}

\begin{document}

\intestazione{TEEA --- 4MAT \per{} Argomento 2 \per{} FORMULARIO}
             {L'impedenza}
             {Prof. Pasquale Naclerio \per{} Manutenzione e Assistenza Tecnica}

\begin{multicols}{2}

\begin{blocco}{Impedenza}
\F{\bar{Z} = |Z| \angle \varphi}
\F{|Z| = \frac{|V|}{|I|} \qquad \varphi = \varphi_V - \varphi_I}
\leg{$|Z|$ in \ohm{} \quad $\varphi$ in gradi, tra $-180^\circ$ e $+180^\circ$ \quad moduli = valori efficaci}
\end{blocco}

\begin{blocco}{Legge di Ohm in alternata}
\F{|I| = \frac{|V|}{|Z|} \qquad \varphi_I = \varphi_V - \varphi}
\F{|V| = |Z| \cdot |I| \qquad \varphi_V = \varphi_I + \varphi}
\end{blocco}

\begin{blocco}{Tipo di carico}
{\small
\begin{tabularx}{\linewidth}{@{}l X l@{}}
\toprule
$\varphi$ & carico & corrente\\
\midrule
$0^\circ$ & resistivo puro & in fase\\
$0^\circ < \varphi < 90^\circ$ & ohmico-induttivo & in ritardo\\
$+90^\circ$ & induttivo puro & ritardo $90^\circ$\\
$-90^\circ < \varphi < 0^\circ$ & ohmico-capacitivo & in anticipo\\
$-90^\circ$ & capacitivo puro & anticipo $90^\circ$\\
\bottomrule
\end{tabularx}}
\leg{con R, L, C: $|\varphi| \le 90^\circ$}
\end{blocco}

\begin{blocco}{Componenti ideali}
\F{\bar{Z}_R = R \angle 0^\circ}
\F{X_L = 2\pi f L \qquad \bar{Z}_L = X_L \angle {+90^\circ}}
\F{X_C = \frac{1}{2\pi f C} \qquad \bar{Z}_C = X_C \angle {-90^\circ}}
\leg{$2\pi f = \omega$ \quad $L$ in H, $C$ in F, $f$ in Hz}
\leg{in continua: L = cortocircuito, C = circuito aperto}
\leg{$1\U{mH} = 10^{-3}\U{H}$ \quad $1\,\mu\mathrm{F} = 10^{-6}\U{F}$ \quad $1\U{nF} = 10^{-9}\U{F}$ \quad $1\U{kHz} = 10^{3}\U{Hz}$}
\end{blocco}

\columnbreak

\begin{blocco}{Formule inverse}
\F{f = \frac{X_L}{2\pi L} \qquad L = \frac{X_L}{2\pi f}}
\F{f = \frac{1}{2\pi C X_C} \qquad C = \frac{1}{2\pi f X_C}}
\end{blocco}

\begin{blocco}{Circuiti in serie (stessa corrente)}
\F{X = X_L - X_C}
\F{|Z| = \sqrt{R^2 + X^2} \qquad \varphi = \arctan\frac{X}{R}}
\F{V_R = R\,|I| \qquad V_X = |X|\,|I|}
\F{|V| = \sqrt{V_R^2 + V_X^2}}
\leg{$X > 0$: induttivo \quad $X < 0$: capacitivo \quad RL: $X = X_L$ \quad RC: $X = -X_C$}
\end{blocco}

\begin{blocco}{Circuiti in parallelo (stessa tensione)}
\F{I_R = \frac{|V|}{R} \quad I_L = \frac{|V|}{X_L} \quad I_C = \frac{|V|}{X_C}}
\F{|I| = \sqrt{I_R^2 + (I_C - I_L)^2}}
\F{\varphi_I = \arctan\frac{I_C - I_L}{I_R}}
\F{|Z| = \frac{|V|}{|I|} \qquad \varphi = -\varphi_I}
\leg{con $\bar{V}$ a $0^\circ$ \quad $I_R$ a $0^\circ$, $I_L$ a $-90^\circ$, $I_C$ a $+90^\circ$}
\end{blocco}

\begin{blocco}{Risonanza ($X_L = X_C$)}
\F{f_0 = \frac{1}{2\pi\sqrt{L C}}}
\leg{serie: $\bar{Z} = R \angle 0^\circ$, $|Z|$ minima, corrente massima e in fase}
\leg{parallelo: $\bar{Z} = R \angle 0^\circ$, $|Z|$ massima, corrente minima}
\end{blocco}



\end{multicols}

\end{document}
